\section*{Capítulo \ref{ch:g8:lenses-images} --- As lentes e as imagens}

\begin{solution}{exo:g8:lenses-images:1}
Barriga no meio: convergente --- ela entorta os raios de um \omterm{def:g6:light-propagation:ray}{feixe}
paralelo um na direção do outro, recolhendo-os no \omterm{def:g8:lenses-images:focus}{foco}. Cova no meio:
divergente --- ela abre o \omterm{def:g6:light-propagation:ray}{feixe} em leque.
\end{solution}

\begin{solution}{exo:g8:lenses-images:2}
O \omterm{def:g8:lenses-images:focus}{foco} é onde um \omterm{def:g6:light-propagation:ray}{feixe} paralelo se recolhe depois da \omterm{def:g8:lenses-images:families}{lente}; a
\omterm{def:g8:lenses-images:focus}{distância focal} $f$ é a distância da \omterm{def:g8:lenses-images:families}{lente} ao \omterm{def:g8:lenses-images:focus}{foco}. Um $f$ curto quer
dizer curvatura forte e entortada com força.
\end{solution}

\begin{solution}{exo:g8:lenses-images:3}
Jogue a \omterm{prop:g5:mirrors-reflection:image}{imagem} da cena distante numa parede, deixe-a nítida movendo a
\omterm{def:g8:lenses-images:families}{lente} e meça da \omterm{def:g8:lenses-images:families}{lente} à parede: essa distância é $f$. A cena precisa
ser distante para que os raios dela cheguem praticamente paralelos
--- e raios paralelos, por definição do \omterm{def:g8:lenses-images:focus}{foco}, se recolhem exatamente
em $f$.
\end{solution}

\begin{solution}{exo:g8:lenses-images:4}
Nunca concentre \omterm{def:g1:light-and-shadows:source}{luz} do Sol perto de nada inflamável --- o pontinho do
\omterm{def:g8:lenses-images:focus}{foco} carboniza e incendeia em segundos. Nunca olhe para o Sol através
de \omterm{def:g8:lenses-images:families}{lente} nenhuma --- o olho mais a \omterm{def:g8:lenses-images:families}{lente} focalizam a queimadura sobre
a retina, na hora e sem dor.
\end{solution}

\begin{solution}{exo:g8:lenses-images:5}
Uma figura montada de verdade por raios recolhidos, apanhável numa
tela. Assinaturas: invertida e clara --- a \omterm{def:g8:lenses-images:families}{lente} colhe uma \omterm{def:g8:lenses-images:families}{lente}
inteira de \omterm{def:g1:light-and-shadows:source}{luz} por ponto, onde o furinho deixava passar um fiozinho
só.
\end{solution}

\begin{solution}{exo:g8:lenses-images:6}
\omterm{def:g8:lenses-images:families}{Lente}: a objetiva de vidro --- a \omterm{def:g8:lenses-images:families}{lente} mole do olho, atrás da pupila.
Tela: o sensor --- a retina. Focalização: a câmera desliza a \omterm{def:g8:lenses-images:families}{lente}
dela; o olho remodela a \omterm{def:g8:lenses-images:families}{lente} dele, mais gorda para perto, mais fina
para longe.
\end{solution}

\begin{solution}{exo:g8:lenses-images:7}
A miopia focaliza forte demais --- as imagens se montam antes da
retina: uma \omterm{def:g8:lenses-images:families}{lente divergente} desentorta na medida certa. O olho
hipermetrope que envelhece focaliza fraco demais para o trabalho de
perto: uma \omterm{def:g8:lenses-images:families}{lente convergente} empresta a \omterm{def:g2:pushes-and-pulls:force}{força} que falta.
\end{solution}

\begin{solution}{exo:g8:lenses-images:8}
O selo precisa ficar mais perto da \omterm{def:g8:lenses-images:families}{lente} do que a \omterm{def:g8:lenses-images:focus}{distância focal}
dela. O olho desfruta então de um fantasma direito e aumentado ---
raios entortados como se viessem de um selo mais grandioso atrás do
vidro; nada cai em tela nenhuma.
\end{solution}

\begin{solution}{exo:g8:lenses-images:9}
Conforme o objeto se aproxima, a \omterm{prop:g5:mirrors-reflection:image}{imagem} dele recua --- então a câmera
afasta a \omterm{def:g8:lenses-images:families}{lente} do sensor para perseguir a \omterm{prop:g5:mirrors-reflection:image}{imagem} em retirada. O olho,
em vez disso, puxa a \omterm{def:g8:lenses-images:families}{lente} dele mais gorda, entortando com mais \omterm{def:g2:pushes-and-pulls:force}{força}
para trazer a \omterm{prop:g5:mirrors-reflection:image}{imagem} de volta à retina, que é fixa.
\end{solution}

\begin{solution}{exo:g8:lenses-images:10}
Na \omterm{def:g8:lenses-images:focus}{distância focal}: só ali a \omterm{def:g8:lenses-images:families}{lente} aglomera os raios praticamente
paralelos do Sol no pontinho menor e mais feroz. Mais perto ou mais
longe, a \omterm{def:g1:light-and-shadows:source}{luz} se espalha por uma mancha suave demais para incendiar
--- o fogo espera precisamente em $f$.
\end{solution}

\begin{solution}{exo:g8:lenses-images:11}
Os óculos de leitura são convergentes: eles recolhem a \omterm{def:g1:light-and-shadows:source}{luz} do Sol num
\omterm{def:g8:lenses-images:focus}{foco} que queima e ampliam o texto de perto --- os dois talentos são
exclusividade das convergentes. As \omterm{def:g8:lenses-images:families}{lentes divergentes} do neto míope
abrem todo \omterm{def:g6:light-propagation:ray}{feixe}: sem \omterm{def:g8:lenses-images:focus}{foco}, sem pontinho que queima, e o texto visto
por elas encolhe em vez de crescer.
\end{solution}

\begin{solution}{exo:g8:lenses-images:12}
Telescópio, estágio um: a grande \omterm{def:g8:lenses-images:families}{lente} recebe os raios praticamente
paralelos de uma \omterm{def:g4:solar-system:star}{estrela} e os recolhe numa \omterm{prop:g8:lenses-images:realimage}{imagem real} pequena perto
do \omterm{def:g8:lenses-images:focus}{foco} dela; o estágio dois amplia essa \omterm{prop:g5:mirrors-reflection:image}{imagem}. A \omterm{def:g1:light-and-shadows:source}{luz} fraca das
\omterm{def:g4:solar-system:star}{estrelas} é preciosa --- a \emph{largura} da \omterm{def:g8:lenses-images:families}{lente} da frente decide
quanto dela é colhido por \omterm{def:g4:solar-system:star}{estrela}: a abertura é tudo. Microscópio,
estágio um: uma \omterm{def:g8:lenses-images:families}{lente} muito curva, de $f$ curto, pendurada bem em
cima da gotinha joga uma \omterm{prop:g8:lenses-images:realimage}{imagem real} muito aumentada; o estágio dois
amplia de novo. O objeto dele é claro e está perto --- o que o
estágio um precisa não é largura, é poder de entortar. Mesma
arquitetura, apetites opostos.
\end{solution}

\begin{solution}{pb:g8:lenses-images:1}
\textbf{1.} Tato: meio com barriga, convergente; meio com cova,
divergente. Teste da janela: as \omterm{def:g8:lenses-images:families}{lentes convergentes} jogam uma janela
de cabeça para baixo na parede a alguma distância; as divergentes
nunca jogam.
\textbf{2.} $f = \qty{25}{cm}$; convergente.
\textbf{3.} Divergente --- \omterm{prop:g8:lenses-images:realimage}{imagem real} nenhuma a distância nenhuma, e
a vista encolhida através dela sela o veredito.
\textbf{4.} C: a \omterm{def:g8:lenses-images:focus}{distância focal} mais curta (\qty{10}{cm}) marca a
\omterm{def:g8:lenses-images:families}{lente} mais gorda e que entorta com mais \omterm{def:g2:pushes-and-pulls:force}{força}.
\textbf{5.} Borrão de longe com visão de perto tranquila: miopia ---
o olho entorta forte demais. Receite a família divergente.
\textbf{6.} Leitura com o braço esticado: o olho hipermetrope que
envelhece e focaliza fraco. Receite \omterm{def:g8:lenses-images:families}{lentes} de leitura convergentes.
\textbf{7.} As \omterm{def:g8:lenses-images:families}{lentes} de leitura acrescentam \omterm{def:g2:pushes-and-pulls:force}{força} de entortar
orçada para páginas de perto; apontadas para as montanhas, cujos
raios chegam paralelos e não precisam de ajuda, a \omterm{def:g2:pushes-and-pulls:force}{força} acrescentada
entorta demais --- as imagens se montam antes da retina, e os picos
borram.
\textbf{8.} Acrescentar correção a um olho bem focalizado o
desequilibra: a \omterm{def:g8:lenses-images:families}{lente} de teste entorta uma \omterm{def:g1:light-and-shadows:source}{luz} que o olho da criança
precisa então combater, montando imagens fora da retina --- a tabela
nadando é o protesto do olho saudável.
\textbf{9.} A \omterm{def:g8:lenses-images:families}{lente} da frente recolhe os raios paralelos da \omterm{def:g4:solar-system:star}{estrela}
numa \omterm{prop:g8:lenses-images:realimage}{imagem real} pequena perto do \omterm{def:g8:lenses-images:focus}{foco} dela; a \omterm{def:g8:lenses-images:families}{lente} do olho é uma
lupa pendurada dentro da própria \omterm{def:g8:lenses-images:focus}{distância focal} em relação a essa
\omterm{prop:g5:mirrors-reflection:image}{imagem}, ampliando-a dentro do olho de quem observa.
\textbf{10.} Cada \omterm{def:g4:solar-system:star}{estrela} entrega apenas um sussurro de \omterm{def:g1:light-and-shadows:source}{luz} por
\omterm{def:g2:measuring-length:units}{centímetro} quadrado; a \omterm{def:g8:lenses-images:families}{lente} larga como um pires colhe \omterm{def:g2:measuring-length:units}{centímetros} às
centenas, afunilando um prato inteiro de sussurro num único ponto
brilhante. A largura é o apetite do astrônomo.
\textbf{11.} A lupa do tamanho de uma moeda, por mais curva que seja,
colheria uma moeda de \omterm{def:g1:light-and-shadows:source}{luz} de \omterm{def:g4:solar-system:star}{estrela} onde é preciso um pires:
imagens nítidas, talvez, mas famintas --- as \omterm{def:g4:solar-system:star}{estrelas} fracas somem.
Ganho: nada que a \omterm{def:g8:lenses-images:families}{lente} do olho já não desse. Veredito: substituição
nenhuma; encomende vidro de verdade.
\textbf{12.} Por exemplo: ``Tudo nesta prateleira faz dois verbos:
\emph{recolher} \omterm{def:g1:light-and-shadows:source}{luz} e depois \emph{entregá-la} onde o olho precisa
--- e a caixa de ferramentas óptica inteira da civilização, dos
óculos aos observatórios, são esses dois verbos em tamanhos
diferentes.''
\end{solution}
